% ---
% id: transmitter/decoy/rem:no-intensity-split
% kind: rem
% title: "Polarisation flags shared across intensities"
% summary: >-
%   A flag family split by intensity would still admit a label-independent
%   tagging channel, but it would drop the relation between one photon number's
%   behaviour at different intensities, which is the decoy content of the
%   declaration. The emission therefore refuses it with a named refusal.
% prominence: standard
% status: accepted
% uses: [def:tagged-source]
% mentions: [prop:tagging, req:shared-flags]
% assumes: []
% refs: []
% provenance: newly-derived
% lean:
% notes: "Deliberate wording: the remark does not call splitting a security error, since a split declaration still admits a channel of the shape prop:tagging constructs and the source identity does not fail; it names the mechanism of the loss and its direction, and asserts no comparison of two key lengths."
% ---

\begin{remark}
\label{rem:no-intensity-split}
Definition~\ref{def:tagged-source}\,(T3) declares \emph{one} polarisation flag
$f_{\mathsf{n},a}$ per photon number $N_{\mathrm{ent}}<\mathsf{n}\le
N_{\mathrm{src}}$ and per polarisation, shared by every intensity and both
round types. A declarer might instead declare a finer family
$f_{\mathsf{n},a,\mu}$, one flag per intensity as well; this remark says what
that would and would not be.

\emph{The source identity would still hold.} The finer declaration admits a
channel of the shape Proposition~\ref{prop:tagging} constructs, and with it the
source identity, for this link \eqref{eq:tagid}: send $f_{\mathsf{n},a,\mu}$ to
$\dyad{\mathsf{n},a}$, which is independent of $\mu$, and leave the rest as it
is. So splitting is not excluded by the source identity and is not a security
error of the kind this submodule is about.

\emph{The decoy content would be lost.} With one shared flag, the letters
$(a,\mu,z)$ and $(a,\mu',z')$ of the same polarisation put their weight on the
same one-dimensional sector of $\Avirt$, so one sector of the declaration
carries the whole of that photon number's weight across intensities and round
types. With one flag per intensity they are different sectors, and nothing in
the declaration relates what happens in the sector of $\mu$ to what happens in
the sector of $\mu'$: a tagged photon number of one intensity and the same
photon number of another intensity become unrelated parts of the declaration.
That relation is the decoy content of the declaration, the reason for declaring
intensities at all being that the same photon number behaves the same way
whichever intensity produced it; so a split declaration is a declaration of a
different protocol, one with intensities and without decoy. The same applies to
the round type: the shared flag is why a tagged photon number of a test round
and of a generation round are one sector, which is not an assumption added here
but a consequence of declaring one flag.

\emph{The direction of the loss.} Splitting removes a relation from the
declaration and adds none, so a consumer of a split declaration is told less
about the source than a consumer of the shared one. The loss would not show as
an error but as a rate that does not improve when an intensity is added, which
is why the emission, the step that outputs the declaration, refuses a split
family with a named refusal (Requirement~\ref{req:shared-flags}). No claim of
this repository compares the key lengths of two declarations, and this remark
asserts no such comparison: it asserts which relation is present in which
declaration. The reverse, a coarser flag than
Definition~\ref{def:tagged-source}\,(T3) declares (for instance one flag for
two polarisations of the same photon number), is excluded for a stricter
reason: it is not a \tref{tagged source declaration} at all. Which vectors the flags
are and which letters put weight on them is fixed by
Definition~\ref{def:tagged-source}\,(T4), and Proposition~\ref{prop:tagging}
proves the source identity for those declarations only.
\end{remark}
